% !TEX root = ../mainfile.tex
% !TEX encoding = UTF-8 Unicode
% !TEX spellcheck = en_US

\chapter{Appendix}
\label{ch:appendix_a}

\section{Currency Universe and Classification}
\label{sec:app_classification}

Table~\ref{tab:app_ccy} lists the 32 currencies in the panel with their carry-sorted bucket
(low/middle/high forward-discount tercile, pegged and managed currencies excluded from the
sort) and their oil-trade role. The literature-standard robustness baskets used in
Chapter~\ref{ch:results} are the yen, Swiss franc and euro (safe) against the Turkish
lira, Brazilian real, South African rand, Mexican peso, Colombian peso and Indonesian rupiah
(risky).

\begin{table}[htbp]
\centering
\begin{threeparttable}
\caption{Currency classification}
\label{tab:app_ccy}
\small
\begin{tabular*}{\textwidth}{@{\extracolsep{\fill}} ll l l @{}}
\toprule
Code & Currency & Carry bucket & Oil role \\
\midrule
AUD & Australian dollar   & middle            & --- \\
BRL & Brazilian real      & no 1M forward     & exporter \\
CAD & Canadian dollar     & middle            & exporter \\
CHF & Swiss franc         & low (safe)        & --- \\
CLP & Chilean peso        & no 1M forward     & --- \\
CNY & Chinese yuan        & pegged/managed    & --- \\
COP & Colombian peso      & no 1M forward     & exporter \\
CZK & Czech koruna        & middle            & --- \\
DKK & Danish krone        & pegged (EUR)      & --- \\
EUR & Euro                & low (safe)        & --- \\
GBP & Pound sterling      & middle            & --- \\
HKD & Hong Kong dollar    & pegged (USD)      & --- \\
HUF & Hungarian forint    & high (risky)      & --- \\
IDR & Indonesian rupiah   & high (risky)      & --- \\
ILS & Israeli shekel      & low (safe)        & --- \\
INR & Indian rupee        & high (risky)      & importer \\
JPY & Japanese yen        & low (safe)        & importer \\
KRW & South Korean won    & middle            & importer \\
KWD & Kuwaiti dinar       & pegged            & exporter \\
MXN & Mexican peso        & high (risky)      & exporter \\
MYR & Malaysian ringgit   & middle            & --- \\
NOK & Norwegian krone     & middle            & exporter \\
NZD & New Zealand dollar  & middle            & --- \\
PEN & Peruvian sol        & no 1M forward     & --- \\
PLN & Polish zloty        & high (risky)      & --- \\
SAR & Saudi riyal         & pegged            & exporter \\
SEK & Swedish krona       & low (safe)        & --- \\
SGD & Singapore dollar    & middle            & --- \\
THB & Thai baht           & low (safe)        & importer \\
TRY & Turkish lira        & high (risky)      & importer \\
TWD & Taiwan dollar       & low (safe)        & importer \\
ZAR & South African rand  & high (risky)      & --- \\
\bottomrule
\end{tabular*}
\begin{tablenotes}[flushleft]\footnotesize
\item \textit{Notes.} Carry buckets are terciles of the average one-month forward discount
(Section~\ref{sec:data_portfolios}). Pegged/managed currencies and the four currencies
without a clean one-month forward (BRL, CLP, COP, PEN) are excluded from the carry sort. Oil roles use observed crude oil trade over 2019--2024 (Section~\ref{sec:data_portfolios}). Partner reports retrieved from the UN Comtrade API on 15~September 2026 supplement 31 of 33 missing country-flow-year observations across India, Japan, Kuwait, Mexico, Taiwan and Saudi Arabia. Kuwait's imports in 2019 and 2023 remain unobserved, not zero. For Türkiye, partner reports imply imports of 28.58 billion dollars and exports of 9.56 billion dollars, an observed import surplus of 19.02 billion dollars. Its import mirror contains no Iraqi reports for 2019--2024 or Russian reports for 2022--2024. Complete partner coverage cannot be established, and mirror values may differ from national totals. All observed balances support the stated roles, conditional on the remaining coverage gaps.
\end{tablenotes}
\end{threeparttable}
\end{table}

\section{Robustness}
\label{sec:app_robustness}

\paragraph{Market model.} Alongside the constant-mean specification, the event study computes
market-model abnormal returns, using the equal-weighted dollar basket as the market return
(Code/\texttt{03\_event\_study.py}, method \texttt{market\_model}).
Table~\ref{tab:market_model} reports the basket-level results: the safe-haven ordering under both treatment shocks is unchanged relative to the constant-mean results in Chapter~\ref{ch:results}.

% Auto-generated by Code/event_study/make_mm_table.py — do not edit by hand.
\begin{table}[htbp]
\centering
\begin{threeparttable}
\caption[Market-model CARs (robustness)]{Market-model CARs of the currency baskets}
\label{tab:market_model}
\footnotesize
\setlength{\tabcolsep}{2pt}
\begin{tabular*}{\textwidth}{@{\extracolsep{\fill}} lcccccccc @{}}
\toprule
 & \multicolumn{2}{c}{Ukraine} & \multicolumn{2}{c}{Liberation Day} & \multicolumn{2}{c}{12-day war} & \multicolumn{2}{c}{Hormuz} \\
\cmidrule(lr){2-3}\cmidrule(lr){4-5}\cmidrule(lr){6-7}\cmidrule(lr){8-9}
 & (0,5) & (0,20) & (0,5) & (0,20) & (0,5) & (0,20) & (0,5) & (0,20) \\
\midrule
Safe currencies & $-0.31\phantom{^{***}}$ & $-1.73^{**}\phantom{^{*}}$ & $+1.31^{***}$ & $+1.17^{*}\phantom{^{**}}$ & $-0.13\phantom{^{***}}$ & $+0.21\phantom{^{***}}$ & $+0.67^{**}\phantom{^{*}}$ & $+0.67\phantom{^{***}}$ \\
Risky currencies & $-0.89\phantom{^{***}}$ & $-0.21\phantom{^{***}}$ & $-0.33\phantom{^{***}}$ & $-0.48\phantom{^{***}}$ & $-0.26\phantom{^{***}}$ & $+0.31\phantom{^{***}}$ & $-0.73^{***}$ & $-1.20^{***}$ \\
Oil exporters & $+1.26^{***}$ & $+2.29^{***}$ & $-0.62\phantom{^{***}}$ & $-0.77\phantom{^{***}}$ & $+0.34\phantom{^{***}}$ & $+0.10\phantom{^{***}}$ & $+0.61\phantom{^{***}}$ & $+1.23^{*}\phantom{^{**}}$ \\
Oil importers & $+0.51\phantom{^{***}}$ & $-1.97\phantom{^{***}}$ & $+1.00^{***}$ & $+0.95\phantom{^{***}}$ & $-0.49\phantom{^{***}}$ & $-0.82\phantom{^{***}}$ & $-0.005\phantom{^{***}}$ & $-0.02\phantom{^{***}}$ \\
Carry (risky$-$safe) & $-0.57\phantom{^{***}}$ & $+1.52\phantom{^{***}}$ & $-1.63^{***}$ & $-1.65\phantom{^{***}}$ & $-0.13\phantom{^{***}}$ & $+0.10\phantom{^{***}}$ & $-1.40^{***}$ & $-1.87^{**}\phantom{^{*}}$ \\
\bottomrule
\end{tabular*}
\begin{tablenotes}[flushleft]\footnotesize
\item \textit{Notes.} Cumulative abnormal returns of the currency baskets in percent.
Abnormal returns are residuals from a regression of each basket on the equal-weighted
dollar basket estimated over the same $[-140,-21]$ window (Section~\ref{sec:method_event}).
The safe-haven ordering under both treatment shocks is unchanged.
$^{*}p<0.10$, $^{**}p<0.05$, $^{***}p<0.01$.
\end{tablenotes}
\end{threeparttable}
\end{table}


\paragraph{Benchmark and persistence figures.} Figure~\ref{fig:car_ukraine_app} shows the
full portfolio response around the external Ukraine benchmark at the identification window.
In the main text the Ukraine dollar path enters only through the
cross-event differences of Figure~\ref{fig:cross_event_usd}. Figure~\ref{fig:car_liberation_w50_app} extends the Liberation Day window to $+50$ days, plotted from $-10$ to $+50$. The tariff drift deepens rather than reverts, the mirror image of the transitory Hormuz premium in Figure~\ref{fig:car_hormuz_w50}.

\begin{figure}[htbp]
  \centering
  \begin{subfigure}{0.48\linewidth}
    \includegraphics[width=\linewidth]{fig_car_ukraine_usd_half}
  \end{subfigure}\hfill
  \begin{subfigure}{0.48\linewidth}
    \includegraphics[width=\linewidth]{fig_car_ukraine_brent_half}
  \end{subfigure}
  \\[0.8em]
  \begin{subfigure}{0.48\linewidth}
    \includegraphics[width=\linewidth]{fig_car_ukraine_safe_half}
  \end{subfigure}\hfill
  \begin{subfigure}{0.48\linewidth}
    \includegraphics[width=\linewidth]{fig_car_ukraine_risky_half}
  \end{subfigure}
  \\[0.8em]
  \begin{subfigure}{0.48\linewidth}
    \includegraphics[width=\linewidth]{fig_car_ukraine_oilexp_half}
  \end{subfigure}\hfill
  \begin{subfigure}{0.48\linewidth}
    \includegraphics[width=\linewidth]{fig_car_ukraine_oilimp_half}
  \end{subfigure}
  \caption[Ukraine invasion: portfolio CARs]{Ukraine invasion (24~February 2022): cumulative abnormal returns over $-5$ to $+20$ trading days. Construction as in Figure~\ref{fig:car_liberation_day}.}
  \label{fig:car_ukraine_app}
\end{figure}

\begin{figure}[htbp]
  \centering
  \begin{subfigure}{0.48\linewidth}
    \includegraphics[width=\linewidth]{fig_w50_liberation_day_usd_half}
  \end{subfigure}\hfill
  \begin{subfigure}{0.48\linewidth}
    \includegraphics[width=\linewidth]{fig_w50_liberation_day_brent_half}
  \end{subfigure}
  \\[0.8em]
  \begin{subfigure}{0.48\linewidth}
    \includegraphics[width=\linewidth]{fig_w50_liberation_day_oilexp_half}
  \end{subfigure}\hfill
  \begin{subfigure}{0.48\linewidth}
    \includegraphics[width=\linewidth]{fig_w50_liberation_day_oilimp_half}
  \end{subfigure}
  \caption[Liberation Day: CARs over $+50$ days]{Liberation Day, $-10$ to $+50$ trading days, estimation window $[-170,-51]$. Construction as in Figure~\ref{fig:car_liberation_day}. The post-announcement dollar decline deepens over the horizon. The 9~April tariff pause (marked) interrupts but does not reverse it.}
  \label{fig:car_liberation_w50_app}
\end{figure}

\paragraph{Regression robustness.} The safe-haven regression of
Section~\ref{sec:results_regression} uses Newey--West (HAC) standard errors, so its inference is
robust to the autocorrelation and volatility clustering of daily returns. This is the main reason a
simple least-squares interaction is preferred here to a conditional-correlation model. The
crisis windows are the same event-day-to-$+20$ windows used in the event study, so the
regression and the event study rest on a common, transparent definition of each episode.
